\section{Sparse moment information}\label{sec:intro}

Can \(O(n)\) prescribed pairwise moments enforce a constant loss in the
best expected coverage, even when almost all pairs remain unrestricted?
Consider a coupling that selects an entire block: it can cover every
feature while keeping each vertex's marginal small. Testing pairs within
blocks restricts precisely the collisions that make this coupling valuable.
We quantify the resulting loss, identifying a sharp hierarchy indexed by
a coloring bound. The information consists of exact probability constraints;
statistical estimation of those probabilities is a separate question.

\paragraph{Model.}
For a finite ground set \(V\), a simple graph \(G=(V,E)\), and
\(x\in[0,1]^V\), let \(\D(x)\) be all laws of \(S\subseteq V\)
with \(\Prb(i\in S)=x_i\). Define
\begin{align*}
\D_G^{\eq}(x)&=\{D\in\D(x):\Prb_D(i,j\in S)=x_ix_j\ (ij\in E)\},\\
\D_{G,\delta}^{\up}(x)&=\{D\in\D(x):\Prb_D(i,j\in S)\leq(1+\delta)x_ix_j\ (ij\in E)\},
\end{align*}
where \(\delta\geq0\) is finite. The graph specifies tested pairs:
nonedges impose no restrictions, unlike a dependency-graph convention.
In particular \(\D_G^{\eq}(x)\subseteq\D_{G,0}^{\up}(x)\), since
upper constraints also permit negative correlations.
A nonnegative weighted coverage function has the form
\[
 f(S)=\sum_{H\in\mathcal H}w_H\mathbf1\{S\cap H\ne\varnothing\},
 \qquad w_H\geq0,
\]
with a finite feature family and finite weights. Set
\[
 f^+(x)=\max_{D\in\D(x)}\E_D f(S),\quad
 f_G^{\eq}(x)=\max_{D\in\D_G^{\eq}(x)}\E_D f(S),\quad
 f_{G,\delta}^{\up}(x)=\max_{D\in\D_{G,\delta}^{\up}(x)}\E_D f(S).
\]
All maxima exist: their feasible sets are nonempty closed subsets of a finite
probability simplex, containing the product law. We consider ratios with
\(f^+(x)>0\); a positive-weight feature then contains a positive-marginal
vertex, so the product expectation and every denominator are positive.

For a class \(\mathcal A\) of finite graphs, write
\[
 \Gamma^{\eq}(\mathcal A)=\sup_{G\in\mathcal A,\,x,f:\,f^+(x)>0}
 \frac{f^+(x)}{f_G^{\eq}(x)},\qquad
 \Gamma_\delta^{\up}(\mathcal A)=\sup_{G\in\mathcal A,\,x,f:\,f^+(x)>0}
 \frac{f^+(x)}{f_{G,\delta}^{\up}(x)}.
\]
The ground set, marginals and finite feature family vary in these suprema.
Let \(\C_k=\{G:\chi(G)\leq k\}\), \(\B_k=\{G:\Delta(G)\leq k-1\}\),
and \(\alpha_k=1-(1-1/k)^k\).

\begin{theorem}\label{thm:main}
For every integer \(k\geq1\) and every fixed finite \(\delta\geq0\),
\[
 \Gamma^{\eq}(\C_k)=\Gamma^{\eq}(\B_k)
 =\Gamma_\delta^{\up}(\C_k)=\Gamma_\delta^{\up}(\B_k)=\alpha_k^{-1}.
\]
\end{theorem}

These are class suprema, approached by a sequence of finite witnesses.
For \(k=2\), a matching on \(2m\) vertices suffices: its specified witness
has exact equality-model ratio \(4m/(3m+1)\), tending to \(4/3\)
(Theorem~\ref{thm:matching}). For fixed \(k\), the witnesses use
\(n(k-1)/2\) tested pairs. Table~\ref{tab:constants} records the first
constants and the finite-error coefficient.

\begin{table}[!hbp]
\centering
\caption{Exact constants. The witness's multiplicative finite upper bound is
\(\min\{1,\alpha_k+c_k(1+\delta)/m\}\), with \(c_k=(k-1)/(2k)\).}
\label{tab:constants}
\begin{tabular}{c|ccccc}
\toprule
\(k\)&2&3&4&5&6\\
\midrule
\(\alpha_k\)&\(3/4\)&\(19/27\)&\(175/256\)&\(2101/3125\)&\(31031/46656\)\\
\(\alpha_k^{-1}\)&\(4/3\)&\(27/19\)&\(256/175\)&\(3125/2101\)&\(46656/31031\)\\
\(c_k\)&\(1/4\)&\(1/3\)&\(3/8\)&\(2/5\)&\(5/12\)\\
\bottomrule
\end{tabular}
\end{table}

\paragraph{Two complementary ingredients.}
A coloring converts any marginal-preserving sampler into a feasible law,
retaining an \(\alpha_k\) fraction of its coverage. This gives the class
upper bound. For sharpness we use a known Cartesian coverage family on
disjoint cliques. The obstacle is to control \emph{every} coupling after
cross-block tests have been removed: deleting tests enlarges the denominator.
Our pointwise certificate uses only total occupancy and within-block
collisions, whose expectations are fixed or bounded by the sparse tests.
It also yields error tolerance, exact matching examples and a finite
independence-hierarchy bound.

\section{Relation to prior work}\label{sec:related}

Agrawal, Ding, Saberi and Ye~\cite{Agrawal2012} developed correlation-gap
bounds for stochastic optimization, comparing the unrestricted coupling
with the mutually independent expectation. The multilinear extension and
continuous-greedy framework of Calinescu, Chekuri, P\'al and
Vondr\'ak~\cite{Calinescu2011} provide a related foundation in submodular
optimization under matroid constraints. Our denominator instead maximizes
over laws subject to a prescribed collection of moment constraints.

Ramachandra and Natarajan~\cite{RamachandraNatarajan2025} introduced the
pairwise-independent correlation-gap framework. Their three-variable
pairwise-optimized bound is \(4/3\) (Theorem 5); the dimension-dependent
\(27/19\) value discussed in their Section 1.1 uses the \emph{mutually}
independent denominator. These constants concern different feasible sets.
Their later five-element coverage counterexample~\cite{RamachandraNatarajan2026}
gives pairwise ratio at least \(640/479>4/3\), disproving the proposed
universal pairwise \(4/3\) bound while preserving the three-variable result.
Here \(k\) measures graph colorability, and the number of variables is
\(n=mk\).

Ramachandra~\cite[Theorem 2 and Section 4]{Ramachandra2026} gives the
Cartesian coverage family used below, its unrestricted optimum and the
asymptotic all-pairs gap \(e/(e-1)\), also noting the \(t\)-wise consequence.
Our certificate keeps the block size \(k\) fixed and allows arbitrary
cross-block dependence. It proves the sparse tightness needed for
Theorem~\ref{thm:main}; Corollary~\ref{cor:twise} also recovers the known
limit with an explicit finite bound.

The two-group coverage calculation in Section 3.2 of the full version of
\emph{The Cost of Consistency}, by Paul D\"utting, Federico Fusco, Silvio
Lattanzi, Ashkan Norouzi-Fard, Ola Svensson and Morteza
Zadimoghaddam~\cite{Dutting2025}, gives the \(3/4\) resampling factor.
Lemma~\ref{lem:color} is its elementary \(k\)-group version. Udwani and
Wajc~\cite[Section 3.1]{UdwaniWajc2026} use the Cartesian coverage polynomial,
\(\alpha_k\), and a no-collision specialization in online submodular welfare.
Our collision penalty quantifies the effect of allowing collisions subject
to tested moments. Thus the contribution combines sparse tightness with a
sharp graph-class hierarchy and finite error refinements; the construction,
constant, resampling ingredient and asymptotic full-independence limit have
the preceding origins.

\section{Resampling and a collision certificate}\label{sec:proof}
\subsection{A conversion from a coloring}\label{sec:color}

\begin{lemma}\label{lem:color}
Given a proper \(k\)-coloring of \(G\) and sampler access to a law
\(D\in\D(x)\), \(k\) independent calls produce a law in
\(\D_G^{\eq}(x)\) with expected coverage at least \(\alpha_k\E_D f\),
simultaneously for all nonnegative weighted coverage functions. Consequently
\(f_G^{\eq}(x)\geq\alpha_k f^+(x)\).
\end{lemma}
\begin{proof}
For \(k\geq2\), partition \(V=C_1\sqcup\cdots\sqcup C_k\) into
independent color classes, allowing empty classes. Draw independent
\(S_1,\ldots,S_k\) with law \(D\), and return
\(R=\bigcup_c(S_c\cap C_c)\). Each vertex keeps its marginal; edge endpoints
use different independent calls, so their joint probability is \(x_ix_j\).

For a feature \(H\), put \(a_c=\Prb_D(S\cap H\cap C_c\ne\varnothing)\),
\(s=\sum_c a_c\), and \(b=\Prb_D(S\cap H\ne\varnothing)\).
The union bound gives \(b\leq\min\{1,s\}\). By independence and AM--GM,
\[
 \Prb(R\cap H\ne\varnothing)
 =1-\prod_c(1-a_c)\geq h(s):=1-(1-s/k)^k.
\]
On \([0,k]\), \(h\) is increasing and concave, with \(h(0)=0\) and
\(h(1)=\alpha_k\). Its chord on \([0,1]\) and monotonicity thereafter give
\(h(s)\geq\alpha_k\min\{1,s\}\geq\alpha_k b\).
Multiplying by nonnegative weights and summing proves the claim for the
single common output law. Apply it to a maximizing input law for the last
assertion. When \(k=1\), the graph is edgeless and the unchanged input law
suffices.
\end{proof}

The conversion uses \(O(k|V|)\) membership operations and no feature
queries. Its input is a sampler; obtaining a maximizing law for \(f^+\)
is a distinct optimization problem.

\subsection{The Cartesian witness}\label{sec:construction}
Fix \(k\geq2\) and \(m\geq1\). Take disjoint blocks \(B_1,\ldots,B_m\)
of \(k\) vertices, and test all pairs within each block. For every
\(u\in\{1,\ldots,k\}^m\), give the transversal feature
\(H_u=\{(r,u_r):1\leq r\leq m\}\) weight \(k^{-m}\).
For \(a_r=|S\cap B_r|\), counting missed features gives
\begin{equation}\label{eq:product}
 f(S)=1-\prod_{r=1}^m(1-a_r/k).
\end{equation}
Here \emph{normalized} means \(f(\varnothing)=0\), \(f(V)=1\) and
\(0\leq f(S)\leq1\). The explicit feature family has \(k^m\) members;
\eqref{eq:product} evaluates it from block counts in \(O(n+m)\) arithmetic
operations. Set \(x_i=1/m\). Selecting one whole block uniformly has
these marginals and hits every transversal, so \(f^+(x)=1\).
The tested graph has \(n=mk\) vertices and
\(|E|=m\binom{k}{2}=n(k-1)/2\), with \(\chi(G)=k\) and
\(\Delta(G)=k-1\). Figure~\ref{fig:construction} shows the construction
and the statistic that the moment constraints control.

\begin{figure}[!htbp]
\centering
\begin{tikzpicture}[x=1cm,y=1cm, every node/.style={font=\normalsize},
 vertex/.style={circle,draw=black,fill=white,inner sep=0pt,minimum size=6pt},
 feature/.style={circle,draw=orange!80!black,fill=orange!80,inner sep=0pt,minimum size=8pt}]
\fill[blue!8,rounded corners] (2.0,-.25) rectangle (3.8,1.55);
\foreach \r/\x in {1/0,2/2.2,3/4.4}{
 \coordinate (a\r) at (\x,0);
 \coordinate (b\r) at (\x+1.4,0);
 \coordinate (c\r) at (\x+.7,1.2);
 \draw[thick] (a\r)--(b\r)--(c\r)--cycle;
 \node[vertex] at (a\r) {};
 \node[vertex] at (b\r) {};
 \node[vertex] at (c\r) {};
 \node at (\x+.7,-.45) {\(B_\r\)};
}
\node[feature] at (a1) {};
\node[feature] at (c2) {};
\node[feature] at (b3) {};
\node[anchor=west] at (6.3,1.1) {\(T=\sum_r a_r\)};
\node[anchor=west] at (6.3,.45) {\(C=\sum_r\binom{a_r}{2}\)};
\draw[-{Stealth},thick] (6.35,-.05)--(6.35,-.6);
\node[anchor=west] at (6.65,-.35) {\(f(S)\leq\alpha_k+\beta_k(T-k)+C/k^2\)};
\end{tikzpicture}

\caption{The witness for \(k=m=3\). Solid triangle edges are tested pairs;
there are no cross-block tests. Orange marks identify one transversal
feature, a set of vertices rather than graph edges. Shading selects the
whole block \(B_2\), which hits every transversal. Such draws have
\(C=3\), while equality-tested moments force \(\E C=1\).
The certificate uses only \(T=|S|\) and \(C\).}
\label{fig:construction}
\Description{Three disjoint tested triangles. Orange vertices form one transversal feature. The middle block is shaded to indicate whole-block selection. Total occupancy and within-block collision count enter the displayed pointwise coverage upper bound.}
\end{figure}

\subsection{A pointwise bound and all error models}\label{sec:certificate}
\begin{lemma}\label{lem:dual}
Let \(q=1-1/k\), \(\beta_k=q^k/k\), \(T=|S|\), and
\(C=\sum_r\binom{|S\cap B_r|}{2}\). Every subset satisfies
\begin{equation}\label{eq:dual}
 f(S)\leq\alpha_k+\beta_k(T-k)+C/k^2.
\end{equation}
\end{lemma}
\begin{proof}
For integers \(0\leq a\leq k\), let \(d_a=q^a-(1-a/k)\).
Then \(d_0=0\) and
\[
 d_{a+1}-d_a=(1-q^a)/k,\qquad
 0\leq1-q^a=k^{-1}\sum_{j=0}^{a-1}q^j\leq a/k.
\]
Summing gives \(0\leq d_a\leq\binom a2/k^2\).
Set \(A_r=1-a_r/k\) and \(Q_r=q^{a_r}\), so
\(0\leq A_r\leq Q_r\leq1\). The product identity
\[
 \prod_rQ_r-\prod_rA_r
 =\sum_r(Q_r-A_r)\prod_{s<r}Q_s\prod_{s>r}A_s\leq C/k^2
\]
yields \(f(S)\leq1-q^T+C/k^2\).
For integer \(t\geq0\), \(g(t)=1-q^t\) has decreasing successive
differences \(q^t/k\). Summing differences above \(k\), or backwards
below \(k\), gives
\(g(t)\leq g(k)+\beta_k(t-k)\). Combining the inequalities proves
\eqref{eq:dual}.
\end{proof}

For edge budgets \(e_{ij}\geq0\), define \(f_{G,e}^{\up}(x)\) by
replacing edge equalities with \(\Prb(i,j\in S)\leq x_ix_j+e_{ij}\).
Write \(f_{G,\epsilon}^{\add}\) for the uniform choice \(e_{ij}=\epsilon\).

\begin{theorem}\label{thm:finite}
For the witness in Section~\ref{sec:construction} and any finite
nonnegative edge budgets,
\begin{equation}\label{eq:hetero}
 \alpha_k\leq f_{G,e}^{\up}(x)\leq
 \min\left\{1,\alpha_k+\frac{k-1}{2km}+\frac1{k^2}\sum_{ij\in E}e_{ij}\right\}.
\end{equation}
In particular, for finite \(\epsilon,\delta\geq0\),
\begin{align}
 \alpha_k\leq f_{G,\epsilon}^{\add}(x)&\leq
 \min\{1,\alpha_k+\tfrac{k-1}{2k}(1/m+m\epsilon)\},\label{eq:add}\\
 \alpha_k\leq f_{G,\delta}^{\up}(x)&\leq
 \min\{1,\alpha_k+\tfrac{(1+\delta)(k-1)}{2km}\}.\label{eq:mult}
\end{align}
Also \(\alpha_k\leq f_G^{\eq}(x)\leq
\min\{1,\alpha_k+(k-1)/(2km)\}\).
\end{theorem}
\begin{proof}
Every admissible law has \(\E T=mk/m=k\), and
\[
 \E C=\sum_{ij\in E}\Prb(i,j\in S)
 \leq \frac{m\binom k2}{m^2}+\sum_{ij\in E}e_{ij}.
\]
Average \eqref{eq:dual}; its linear term cancels and its collision coefficient
is nonnegative. Together with \(f\leq1\), this proves the upper bound.
Lemma~\ref{lem:color} and inclusion of equality-feasible laws give the
lower bound. Substituting \(e_{ij}=\epsilon\) or \(\delta/m^2\) gives
\eqref{eq:add} and \eqref{eq:mult}; equality laws are feasible for the
zero-error upper model, giving the final assertion. At \(m=1\), all
vertices are forced present, and every displayed denominator equals one.
\end{proof}

\subsection{Completing the class theorem}\label{sec:classproof}
\begin{proof}[Proof of Theorem~\ref{thm:main}]
For \(k=1\) every graph in either class is edgeless, giving ratio one.
For \(k\geq2\), greedy coloring gives \(\B_k\subseteq\C_k\).
Feasible-set inclusion and Lemma~\ref{lem:color} give the two chains
\begin{align*}
 \Gamma_\delta^{\up}(\B_k)&\leq\Gamma^{\eq}(\B_k)
 \leq\Gamma^{\eq}(\C_k)\leq1/\alpha_k,\\
 \Gamma_\delta^{\up}(\B_k)&\leq\Gamma_\delta^{\up}(\C_k)
 \leq\Gamma^{\eq}(\C_k)\leq1/\alpha_k.
\end{align*}
The witness belongs to \(\B_k\) and, by \eqref{eq:mult}, has upper-model
ratio at least \([\alpha_k+(1+\delta)(k-1)/(2km)]^{-1}\).
Letting \(m\to\infty\) at fixed \(k,\delta\) makes the endpoints of
both chains equal, proving all four identities.
\end{proof}

\section{Exact finite examples and error scales}\label{sec:finite}
\subsection{A matching witness with an exact transition}\label{sec:matching}

\begin{theorem}\label{thm:matching}
For the Cartesian witness with \(k=2\), every integer \(m\geq1\) and
finite \(\delta,\epsilon\geq0\),
\begin{align}
 f_G^{\eq}(x)&=\frac34+\frac1{4m},&
 \frac{f^+(x)}{f_G^{\eq}(x)}&=\frac{4m}{3m+1},\label{eq:matching-eq}\\
 f_{G,\delta}^{\up}(x)&=\frac34+\frac14\min\{1,(1+\delta)/m\},\label{eq:matching-mult}\\
 f_{G,\epsilon}^{\add}(x)&=\frac34+\frac14\min\{1,1/m+m\epsilon\}.\label{eq:matching-add}
\end{align}
For any nonnegative error sequences as \(m\to\infty\), these latter
denominators tend to \(3/4\) if and only if \(\delta_m/m\to0\),
respectively \(m\epsilon_m\to0\). If the respective scaled error tends
to \(c\in[0,\infty]\), the denominator tends to
\(3/4+\min\{1,c\}/4\) and the ratio to \(4/(3+\min\{1,c\})\).
\end{theorem}
\begin{proof}
For \(m\geq2\), type A selects two distinct uniform blocks and one
independent fair vertex in each: its marginals are \(1/m\), edge moments
zero, and coverage \(3/4\). Type B selects one entire uniform block:
it has the same marginals, edge moments \(1/m\), and coverage one.
Using type B with probability \(\theta\) gives moments \(\theta/m\)
and value \(3/4+\theta/4\). For the multiplicative model choose
\(\theta=\min\{1,(1+\delta)/m\}\); for the additive model choose
\(\theta=\min\{1,1/m+m\epsilon\}\). Their values meet the respective
bounds of Theorem~\ref{thm:finite}. For equality take \(\theta=1/m\),
meeting the zero-error upper bound with each edge equality exact.
For \(m=1\), the only feasible set is \(V\), giving one in every formula.
The limit statements follow directly from the formulas and continuity of
\(\min\{1,u\}\) on the extended nonnegative half-line. Necessity follows
because \(\min\{1,u_m\}\to0\) forces \(u_m\to0\) for \(u_m\geq0\).
\end{proof}

\paragraph{General sufficient scales.}
For every fixed \(k\geq2\), Theorem~\ref{thm:finite} implies convergence
of the witness denominator to \(\alpha_k\) and of its ratio to
\(1/\alpha_k\) whenever \(\delta_m=o(m)\), or, in the additive model,
\(\epsilon_m=o(1/m)\). More generally \(\sum_E e_{ij}^{(m)}\to0\)
suffices. The per-edge absolute excess in the multiplicative model is
\(\delta_m/m^2\); there are \(m\binom k2\) tested edges. This explains
the aggregate scale. The necessity assertions of Theorem~\ref{thm:matching}
apply to its specified matching witness and uniform marginals.

\subsection{Five rational optimality certificates}\label{sec:certificates}
A profile is the sorted vector of block occupancies. Averaging a law over
block permutations and independent within-block vertex permutations preserves
coverage and all uniform constraints. The group is transitive on vertices
and on tested edges, so a uniform profile-orbit mixture has marginal
\(\E T/(mk)\) at each vertex and edge moment
\(\E C/[m\binom k2]\). Thus \(\E T=k\) and the relevant constraint on
\(\E C\) suffice for such a mixture to lift to a feasible law.

\begin{proposition}\label{prop:finite}
The four equality optima at \((k,m)=(2,2),(2,3),(3,2),(3,3)\) are
\(7/8,5/6,5/6,65/81\), respectively. At \((3,3)\) the upper-model
optimum with \(\delta=1\) is \(73/81\).
\end{proposition}
\begin{proof}
Table~\ref{tab:certificates} supplies orbit mixtures and pointwise upper
bounds \(f(S)\leq y_0+y_TT+y_CC\). Their masses sum to one,
\(\E T=k\), and \(\E C=k(k-1)/(2m)\) for equality, or twice that
quantity in the last row. Their objectives are the listed values.
For \(k=2\) the pointwise inequality holds as follows: if \(C=0\),
then \(f=1-2^{-T}\leq1/4+T/4\), checked at \(T=0,1,2\) and immediate
for \(T\geq3\); if \(C\geq1\), a full block gives \(f=1\),
\(T\geq2\), and the asserted upper bound is at least one.
For \(k=3,m=2\), the exhaustive profiles in order
\[
 (0,0),(0,1),(0,2),(0,3),(1,1),(1,2),(1,3),(2,2),(2,3),(3,3)
\]
have dual slacks
\(1/3,1/9,0,0,0,0,1/9,1/9,1/3,2/3\).
For \(k=m=3\), Table~\ref{tab:slack} exhausts all 20 profiles;
\eqref{eq:product} depends only on the profile, so these cover all 512
labeled subsets. Every slack is nonnegative. Averaging the dual bound
matches the primal values in each row; \(y_C\geq0\) justifies its use
with the upper constraint. Finite weak duality proves exact optimality.
\end{proof}

\begin{table}[!htbp]
\centering
\caption{Exact primal/dual certificates. Each profile denotes its uniform
subset orbit. All omitted masses are zero; the last row uses \(\delta=1\).}
\label{tab:certificates}
\begin{tabular}{ccclc}
\toprule
\((k,m)\)&Model&Primal profiles: masses&\((y_0,y_T,y_C)\)&Optimum\\
\midrule
\((2,2)\)&eq&\((0,2):1/2;\ (1,1):1/2\)&\((1/4,1/4,1/4)\)&\(7/8\)\\
\((2,3)\)&eq&\((0,0,2):1/3;\ (0,1,1):2/3\)&\((1/4,1/4,1/4)\)&\(5/6\)\\
\((3,2)\)&eq&\((0,3):1/4;\ (1,2):3/4\)&\((1/3,1/9,1/9)\)&\(5/6\)\\
\((3,3)\)&eq&\((0,0,3):1/3;\ (1,1,1):2/3\)&\((8/27,11/81,8/81)\)&\(65/81\)\\
\((3,3)\)&up&\((0,0,3):2/3;\ (1,1,1):1/3\)&\((8/27,11/81,8/81)\)&\(73/81\)\\
\bottomrule
\end{tabular}
\end{table}

\begin{table}[!htbp]
\centering
\caption{Complete pointwise certificate for \(k=m=3\): slack of
\(8/27+(11/81)T+(8/81)C-f(S)\).}
\label{tab:slack}
\begin{tabular}{cc@{\qquad}cc}
\toprule
Profile&Slack&Profile&Slack\\
\midrule
\((0,0,0)\)&\(8/27\)&\((1,1,1)\)&\(0\)\\
\((0,0,1)\)&\(8/81\)&\((1,1,2)\)&\(7/81\)\\
\((0,0,2)\)&\(0\)&\((1,1,3)\)&\(22/81\)\\
\((0,0,3)\)&\(0\)&\((1,2,2)\)&\(20/81\)\\
\((0,1,1)\)&\(1/81\)&\((1,2,3)\)&\(41/81\)\\
\((0,1,2)\)&\(2/81\)&\((1,3,3)\)&\(68/81\)\\
\((0,1,3)\)&\(11/81\)&\((2,2,2)\)&\(4/9\)\\
\((0,2,2)\)&\(4/27\)&\((2,2,3)\)&\(20/27\)\\
\((0,2,3)\)&\(10/27\)&\((2,3,3)\)&\(29/27\)\\
\((0,3,3)\)&\(19/27\)&\((3,3,3)\)&\(38/27\)\\
\bottomrule
\end{tabular}
\end{table}

\paragraph{Reproducibility.}
The accompanying \texttt{finite-certificates.json} stores all five rational
inputs and certificates. The standalone \texttt{verify\_revision.py}
reconstructs uniform orbit laws, checks every marginal and tested pair,
and verifies the objective and dual inequality on every labeled subset
using exact fractions; \texttt{discover\_finite.py} records the separate
rational profile search. These files accompany the source package; \texttt{manuscript/reproduce.py}
provides a read-only launcher.

\section{Consequences and open questions}\label{sec:consequences}
\subsection{A finite full-independence hierarchy}\label{sec:twise}
For the same witness, let \(f_t(x)\) maximize expected coverage over laws
whose indicators are jointly independent on every subfamily of at most
\(t\) vertices, with \(2\leq t\leq n=mk\). Thus \(t=2\) tests all pairs,
and \(t=n\) specifies the product law.

\begin{corollary}\label{cor:twise}
Let \(k,m\geq2\), \(a=1-1/e\), and
\(U=\alpha_k+(k-1)/(2km)\). Uniformly for \(2\leq t\leq mk\),
\begin{align}
 \alpha_m\leq f_t(x)&\leq\min\{1,U\},&
 \frac1U\leq\frac{f^+(x)}{f_t(x)}&\leq\frac1{\alpha_m},\label{eq:twise}\\
 0\leq\frac e{e-1}-\frac{f^+(x)}{f_t(x)}
 &\leq\frac{1/k+1/(2m)}{(1-1/e)^2}.\label{eq:rate}
\end{align}
Consequently any sequences \(k,m\to\infty\), with any admissible choices
of \(t\), give ratio tending to \(e/(e-1)\). On \(k=m=s\geq2\), the
error is at most \(3/[2(1-1/e)^2\sqrt n]\), where \(n=s^2\).
\end{corollary}
\begin{proof}
The product law is feasible and hits each \(m\)-vertex feature with
probability \(1-(1-1/m)^m=\alpha_m\). Every feasible \(t\)-wise law
satisfies the within-clique equalities, so Theorem~\ref{thm:finite} bounds
its objective above by \(\min\{1,U\}\). Since \(f^+=1\), this proves
\eqref{eq:twise}. For \(j\geq2\), \(\log(1-1/j)\leq-1/j\) and
\(\log(1+1/(j-1))\leq1/(j-1)\) imply
\[
 a\leq\alpha_j\leq a+1/(ej)\leq a+1/j.
\]
Thus the ratio is at most \(1/a\), while \(U\geq a\) and
\(U-a\leq1/k+1/(2m)\). Subtracting its lower bound \(1/U\) from
\(1/a\) gives
\((U-a)/(aU)\leq(1/k+1/(2m))/a^2\). This proves the rate and its
square-size specialization.
\end{proof}

The limiting theorem and family are due to Ramachandra~\cite{Ramachandra2026};
the finite estimate follows here from within-block moments alone.
The product expectation is \(\alpha_m\), depending on the number of
vertices per feature: for example \(k=2,m=3,t=6\) gives \(19/27\).

\subsection{Graph structure and optimization questions}
The exact matching formulas isolate how a limited moment budget constrains
a useful unrestricted coupling, and Lemma~\ref{lem:color} converts an
available sampler into a feasible one with a quantified loss. Together
these give finite optimization consequences of the class theorem.
For general fixed graphs, determining the best constant from structure
beyond an integral coloring remains an open direction here. Fractional-color
information and high-chromatic triangle-free graphs are natural tests.
The construction gives a sufficient linear number of tested pairs at fixed
\(k\); a necessary pair budget requires a separate lower bound.
Extending the sharp hierarchy to arbitrary monotone submodular objectives
also remains unresolved.

\paragraph{Computational and AI methodology.}
OpenAI Codex was used for topic selection, proof development, exact
computations, literature checks, drafting and typesetting. A separate
model-context pass checked the proofs and finite certificates.
